\section{Contenerización Docker}

La contenerización de Voctomix 2.0 se articula en dos ficheros: un \texttt{Dockerfile}
que define la imagen base del sistema y un \texttt{docker-compose.yml} que orquesta
el conjunto de servicios que componen la instalación completa.

\subsection{Dockerfile}

La imagen se construye sobre \texttt{ubuntu:22.04} e instala todas las dependencias
del sistema en una sola capa: GStreamer 1.0 con sus plugins base, buenos, malos y
\textit{ugly}, FFmpeg, las \textit{bindings} Python de GObject (\texttt{python3-gi}),
el cliente AMQP (\texttt{pika}) y utilidades de red (\texttt{netcat}, \texttt{iproute2}).
El directorio de trabajo del contenedor es \texttt{/opt/voctomix} y el punto de
entrada es el propio \texttt{voctocore.py}.

Un \textit{healthcheck} basado en \texttt{nc -zw1 localhost 9999} permite a Docker
Compose determinar cuándo el núcleo está listo para aceptar conexiones, lo que evita
condiciones de carrera durante el arranque del sistema.

\subsection{Docker Compose}

El fichero \texttt{docker-compose.yml} define diez servicios organizados en una
red bridge privada (\texttt{voctomix\_net}):

\begin{itemize}
  \item \textbf{voctocore:} núcleo de procesamiento multimedia. Expone los puertos
        de control (9999), salida de stream (15000) y entradas de cámara (10000--10005).
  \item \textbf{cam1--cam4:} fuentes FFmpeg de prueba. Comparten el
        \textit{network namespace} de voctocore para comunicarse por localhost.
  \item \textbf{stream\_blanker:} genera los bucles de pausa y fuera de emisión.
  \item \textbf{rabbitmq:} broker de mensajería AMQP (puertos 5672 y 15672 para
        la consola de administración web).
  \item \textbf{notario:} servicio de telemetría (puerto 8080).
  \item \textbf{audio\_manager:} gestión dinámica de niveles de audio.
\end{itemize}

% --- Figura sugerida ---
% \begin{figure}[H]
%   \centering
%   \includegraphics[width=0.9\textwidth]{figuras/docker_compose_detalle.png}
%   \caption{Diagrama completo de servicios Docker Compose: servicios, red y puertos.}
%   \label{fig:docker_compose_detalle}
% \end{figure}

% --- Tabla sugerida ---
% \begin{table}[H]
%   \centering
%   \caption{Servicios definidos en docker-compose.yml.}
%   \label{tab:servicios_docker}
%   \begin{tabular}{|l|l|c|l|}
%     \hline
%     \textbf{Servicio} & \textbf{Imagen base} & \textbf{Puertos} & \textbf{Función} \\
%     \hline
%     voctocore      & local/voctomix & 9999, 15000, 10000--10005 & Núcleo multimedia \\
%     cam1--cam4     & local/voctomix & (localhost interno)       & Fuentes de prueba \\
%     stream\_blanker& local/voctomix & 17000--17001, 18000       & Bucles pausa/offline \\
%     rabbitmq       & rabbitmq:mgmt  & 5672, 15672               & Broker AMQP \\
%     notario        & local/voctomix & 8080                      & Telemetría REST/AMQP \\
%     audio\_manager & local/voctomix & ---                       & Gestor de audio \\
%     \hline
%   \end{tabular}
% \end{table}
